% concatenated from Z-X-20260924.tex
\documentclass{birkjour}

\usepackage{mathrsfs}
\usepackage{amssymb,amsmath,latexsym}
\usepackage{amsfonts}
\usepackage{booktabs}

% hyperref 只加载一次！不要重复加载 hyperref
\usepackage[colorlinks, citecolor=blue,linkcolor=blue]{hyperref}

%
% THEOREM Environments (Examples)-----------------------------------------
%
\theoremstyle{plain}
\newtheorem{thm}{Theorem}
\newcounter{subthm}[thm]
\renewcommand{\thesubthm}{\thethm.\arabic{subthm}}
\newtheorem{cor}[subthm]{Corollary}
\newtheorem{rem}{Remark}
\newtheorem{lem}{Lemma}
\theoremstyle{definition}
\newtheorem{defn}{Definition}
\theoremstyle{definition}
\newtheorem{theorem}{Theorem}
\renewcommand{\thetheorem}{\Alph{theorem}}
\newtheorem{pro}[thm]{Problem}

\begin{document}


\title[Improved Bohr-Rogosinski radius for holomorphic mappings with values in complex Banach spaces]
{Improved Bohr-Rogosinski radius of holomorphic mappings with values in complex Banach spaces}

\author[L.Y. Zeng, L.P. Xiong]{Luyi Zeng, Liangpeng Xiong$^*$}
\address{School of Mathematics Science, Jiangxi Science and Technology Normal University, Nanchang $330038$
	Jiangxi, People's Republic of China}
\email{lpxiong2016@whu.edu.cn\\
	\\$^*$Correspondence author}

\subjclass{32H02 $\cdot$ 32A30}

\keywords{Bohr-Rogosinski radius $\cdot$ Holomorphic mapping $\cdot$ Homogeneous polynomial expansion $\cdot$ $\text{JB}^*$-triple $\cdot$ Unit polydisc}

\begin{abstract}
In this paper, we extend some Bohr-Rogosinski type inequalities of analytic functions in $\mathbb{C}$ to the cases of holomorphic mappings with values in higher-dimensional complex space. First, we obtain the general Bohr-Rogosinski radius for holomorphic mappings with values in the closure of the unit polydisc in $\mathbb{C} ^n$. Furthermore, we consider the corresponding problems for holomorphic mappings with value in the closure of the unit ball of a JB$^*$-triple. All the radius are optimal.
\end{abstract}

%%% ----------------------------------------------------------------------
\maketitle
%%% ----------------------------------------------------------------------
%\tableofcontents
\section{Introduction} \label{sec1}
Let $\mathcal{B} $ be the class of analytic functions in the unit disk $\mathbb{D}=\{z\in \mathbb{C} : |z|<1\}$  of the form  $f(z)=\sum_{n=0}^{\infty}a_{n}z_{n} $ such that $\left | f(z) \right | < 1$ in $\mathbb{D}$ .
\par
Harald Bohr discovered the famous result related to the family  $\mathcal{B} $ in 1914 \cite{ref9}, the sharp version of the result is the following.

\begin{theorem}\cite{ref9}
	$\mathit{f} \in \mathcal{B}$,
	$f(z)=\sum_{n=0}^{\infty}a_{n}z^{n} \in \mathcal{B}$, then
	\begin{align*}
		\mathit{M}_{f} (r):=\sum_{n=0}^{\infty} \left| a_{n} \right| r^{n} \le 1 \quad
		\text{for} \quad \left| z \right|=r\le \frac{1}{3}
	\end{align*}
and the constant $\frac{1}{3}$ is sharp.
\end{theorem}
\par
The constant $\frac{1}{3}$ and the inequality $\mathit{M}_{f} (r) \le 1$ are called the Bohr radius and the Bohr's inequality for the class $\mathcal{B}$, respectively.
Bohr actually obtained the inequality for $\left| z \right| \le \frac{1}{6}$.
Later, Wiener, Riesz and Schur, independently established the inequality for $\left| z \right| \le \frac{1}{3}$.
By changing the $\left| a_j \right|$ to $\left| a_j \right|^p$, Djakov and Ramanujan \cite{ref18} obtained the Bohr phenomenon from different point of view, they studied the powered Bohr radius $r_p$ with two different settings, $p \in (0,2)$ and $p \in (1,2]$. And this paper also included a discussion of the multidimensional case. A conjecture proposed for powered Bohr radius $r_p$ in \cite{ref18} has been settle affirnatively in \cite{ref29}. Recently, there have been several other refined versions of Bohr radius. For example, the case where $|a_0|$ is replaced by $|a_0|^p$ with $0 < p \leq 2$ can be found in \cite{ref33,ref36}. In this case, the sharp Bohr type radius is $p/(2+p)$. And the case $p=2$ is received a lot of attention and research. In many subsequent articles, the research topic of Bohr radius has attracted the attention of numerous scholars, and the Bohr radius of some power series on the unit disk $\mathbb{D}$ have also been continuously verified (see, e.g.\cite{ref1}, \cite{ref2}, \cite{ref5} \cite{ref7}, \cite{ref20}, \cite{ref27}) and the references therein.
In addition to the Bohr radius, two other radius of analytic functions on the unit disk $\mathbb{D}$ have also recently received attention.


\begin{theorem}(Rogosinski Theorem, \cite{ref37})
	$\mathit{f}(z)=\sum_{n=0}^{\infty}a_{n}z^{n}  \in \mathcal{B}$, then for every N $\ge$ 1, $\mathit{S}_{N} (z):=\sum_{n=0}^{N-1} a_{n} z^{n} $ denotes the partial sums of $\mathit{f}$, then
		\begin{align*}
		\left |\mathit{S}_{N} (z):=\sum_{n=0}^{N-1} a_{n} z^{n} \right | \le 1\quad
		\text{for}\quad \left| z \right| \le \frac{1}{2},
	\end{align*}
	the constant $\frac{1}{2}$ is called the Rogosinski radius for the family $\mathcal{B}$, and the radiu is sharp.
\end{theorem}

 Meanwhile, by combining Bohr's inequality and Rogosinski's inequality, a new inequality has received considerable attention in recent years. Its definition is as follows.

\begin{defn}(Bohr-Rogosinski Theorem, \cite{ref28})
	$\mathit{f}(z)=\sum_{n=0}^{\infty}a_{n}z^{n}  \in \mathcal{B}$, then for every N $\in \mathbb{N}$ , $\mathit{R}_{N}^{f}(z) :=\left| \mathit{f}(z) \right| + \sum_{n=N}^{\infty} \left| a_{n} \right| r^{n}  \le 1 $ denotes the Bohr-Rogosinski sums of $\mathit{f}$, then
	\begin{align*}
		\mathit{R}_{N}^{f}(z) :=\left| f(z) \right| + \sum_{n=N}^{\infty} \left| a_n \right| r^{n}  \le 1\quad
		\text{for}\quad \left| z \right|=r \le R_N,
	\end{align*}
	the inequality is called the Bohr-Rogosinski inequatity for the family $\mathcal{B}$, and the constant $R_N$ is sharp and it is the positive root of the equation
	\begin{equation*}
		2(1+r)r^N-(1-r)^2=0.
	\end{equation*}
\end{defn}
For $\mathit{f}(z)=\sum_{n=0}^{\infty } a_nz^n$, we let $\mathit{f}_0(z):=\mathit f(z)-\mathit f(0)$ and throught the paper, we use the following notations:
\begin{align*}
	\left \| f_0 \right \| ^ 2:=\sum_{n=1}^{\infty } \left | a_n \right | ^2r^{2n}.
\end{align*}

In many studies, they have improved Bohr and Bohr-Rogosinski inequality of the family $\mathcal{B}$ by using $\left \| f_0 \right \| ^ 2$. The improved Bohr-Rogosinski theorem (Theorem D) and improved Bohr theorem (Theorem C) for the family $\mathcal{B}$ were proved by Liu-Liu-Ponnusamy \cite{ref32} and Ponnusamy-Vijayakumar-Wirths \cite{ref35}, respectively.


\begin{theorem}\cite{ref35}
	Suppose that $\mathit{f}(z)=\sum_{n=0}^{\infty} \left| a_{n} \right| r^{n} \in \mathcal{B}$. Then
	\begin{align*}
		\sum_{n=0}^{\infty} \left| a_n \right| r^{n} +(\frac{1}{1+\left | a_0 \right | } +\frac{r}{1-r} )\left \| f_0 \right \| ^ 2 \le 1, \text{for} \left| z \right | =r \le r_0=\frac{1}{2+\left | a_0 \right |}.
	\end{align*}
	and the numbers $\frac{1}{2+\left | a_0 \right |}$ and $\frac{1}{1+\left | a_0 \right |}$ cannot be improved.
\end{theorem}

\begin{theorem}\cite{ref32}
	Suppose that $\mathit{f}(z)=\sum_{n=0}^{\infty} \left| a_{n} \right| r^{n} \in \mathcal{B}$. Then
	\begin{align*}
		\left| \mathit{f}(z) \right|+\sum_{n=N}^{\infty} \left| a_n \right| r^{n} +(\frac{1}{1+\left | a_0 \right | } +\frac{r}{1-r} )\left \| f_0 \right \| ^ 2 \le 1,\\
		\text{for} \left| z \right | =r \le r_0=\frac{2}{3+\left | a_0 \right |+\sqrt{5}(1+a_0)}.
	\end{align*}
	The radius $r_0$ is the best possible and $r_0>\sqrt{5}-2$.
\end{theorem}
Subsequent researchers extended Bohr's inequality from the unit disk to polydisks and more general bounded domains in complex Euclidean spaces \cite{ref3,ref4}, as well as to vector-valued settings \cite{ref34}, which laid the foundation for the systematic study of Bohr-type inequalities.

As commutative bounded symmetric domains, polydisks serve as the most fundamental setting for the multidimensional generalization of the Bohr radius. Dineen and Timoney \cite{ref17} obtained the Bohr radius of the polydisc tends to zero through qualitative analysis. On this basis, Boas and Khavinson \cite{ref8} through systematic research on the Bohr radius $K_n$ on the n-dimensional polydisk and established the two-sided estimate
    \begin{equation*}
       \frac1{3^n}\le K_n\le 2\frac{\log n}{n}.
    \end{equation*}
Obviously, as the dimension $n \to \infty$, the Bohr radius $K_n$ tends to zero. This reveals the fundamental difference between one-dimensional and high-dimensional Bohr phenomena. In subsequent studies, this estimate was continuously improved, see e.g.\cite{ref14,ref16}. And ultimately by Bayart, Pellegrino, and Seoane-Sepúlveda \cite{ref6} determined the exact asymptotics $K_n\sim \frac{\log n}{n}$. Nevertheless, polydisks represent only the reducible, commutative class of bounded symmetric domains, while the four classical families of irreducible bounded symmetric domains in the Cartan classification, together with two exceptional domains, cannot be fully described by the Euclidean dimension alone. A unified algebraic framework is therefore required to treat all symmetric domains simultaneously.

On account of the open unit ball of every $\text{JB}^*$-triple is a bounded symmetric domain; conversely, by Kaup's Riemann mapping theorem, every bounded symmetric domain in a complex Banach space is biholomorphically equivalent to the open unit ball of a $\text{JB}^*$-triple, which is unique up to isometric isomorphism \cite{ref30}. Hence the theory of $\text{JB}^*$-triple systems can provide a standard axiomatic framework that can handle all symmetric domains simultaneously. For the definition and the properties of a $\text{JB}^*$-triple, see e.g.\cite{ref12,ref13}.

Based on this framework, Bohr-type problems have been generalized to the unit balls of $\text{JB}^*$-triples(see e.g.\cite{ref15}, \cite{ref21}, \cite{ref24}, \cite{ref25} for the finite-dimensional case and \cite{ref15}, \cite{ref22}, \cite{ref23}, \cite{ref25}, \cite{ref26} for infinite-dimensional case). Especially in Hamada et al.'s latest research\cite{ref25}, they directly obtained the Bohr–Rogosinski radius for holomorphic mappings on the unit ball of a complex Banach space with values in a higher dimensional complex Banach space and the Bohr–Rogosinski radius for holomorphic mappings with values in the closure of the unit ball of a $\text{JB}^*$-triple. And they also obtained the Bohr-Rogosinski radius for hlomorphic mappings with values in the closure of the unit polydisc of ${\mathbb{C}}^n, n \ge 2$ and in the closure of the unit ball of a $\text{JB}^*$-triples in following theorems.
\begin{theorem}\cite{ref25}
	Let $B_X$ be the unit ball of a complex Banach space $X$. Let $a=(a_1,\dots,a_n)\in\overline{\mathbb{U}}^n$ and let $F$ be a holomorphic mapping from $B_X$ to $\overline{\mathbb{U}}^n$ with
	\begin{equation*}
		F(z)=a+\sum_{s=1}^{\infty}P_s(z),\quad z\in B_X,
	\end{equation*}
	where $P_s(z)=\frac{1}{s!}D^sF(0)(z^s)$. Assume that $|a_j|=\|a\|_{\infty}$ for all $j\in\{1,2,\dots,n\}$. Then for $m,N\in\mathbb{N}$ and $p\in(0,2]$, we have
	\begin{equation*}
		\|F(v(z))\|_{\infty}^p+\sum_{s=N}^{\infty}\|P_s(z)\|_{\infty}\leq 1.
	\end{equation*}
	for $\|z\|=r\leq R_{m,N}^p$, where $v:B_X\to B_X$ is a Schwarz mapping having $z=0$ as a zero of order $m$, and $R_{m,N}^p$ is the unique root in $(0,1)$ of the equation
	\begin{equation*}
		p\left(\frac{1-r^m}{1+r^m}\right)-\frac{2r^N}{1-r}=0.
	\end{equation*}
	The constant $R_{m,N}^p$ cannot be improved.
\end{theorem}
\begin{theorem}\cite{ref25}
	Let $B_X$ be the unit ball of a complex Banach space $X$ and let $\mathbb{B}_Y$ be the unit ball of a $\text{JB}^*$-triple $Y$. Let $F$ be a holomorphic mapping from $B_X$ to $\mathbb{B}_Y$ with
	\[
	F(z) = a + \sum_{s=1}^{\infty} P_s(z),\quad z\in B_X,
	\]
	where $P_s(z) = \frac{1}{s!} D^s F(0)(z^s)$. Assume that $a\in \mathbb{B}_Y$. Then for $m, N\in \mathbb{N}$ and $p\in(0,2]$, we have
	\[
	\|F(v(z))\|^p + \sum_{s=N}^{\infty}\frac{\|D\varphi_a(a)P_s(z)\|}{\|D\varphi_a(a)\|} \le 1
	\]
	for $\|z\| = r \le R_{m,N}^p$, where $v\colon B_X\to B_X$ is a Schwarz mapping having $z=0$ as a zero of order $m$, $\varphi_a\in \operatorname{Aut}(\mathbb{B}_Y)$ such that $\varphi_a(a)=0$ and $R_{m,N}^p$ is the unique root in $(0,1)$ of the equation (2.2). The number $R_{m,N}^p$ cannot be improved.
\end{theorem}
Therefore, a natural question arises: Can the high-dimensional Bohr and Bohr-Rogosinski inequalities given by Theorem E and Theorem F be generalized by using $\left \| f_0 \right \| ^ 2$?
\par
In this paper, we will extend the Bohr and Bohr-Rogosinski type inequalities given by Theorem E and Theorem F to holomorphic mappings with values in the closure of the unit polydisc of ${\mathbb{C}}^n, n \ge 2$ in section 3. And in section 4, we will extend them to holomorphic mappings with value in the closure of the unit ball of a \(JB^*\)-triples. We all obtained the sharp radius.


\section{Preliminaries}\label{Sec:2}

If $X$ and $Y$ be complex Banach spaces with norm $\left\|\bullet\right\|$.
\begin{defn}
	Let $k \in \mathbb{N}$. If there exsits a $k$-linear mapping $u$ from $X^k$ to $Y$ which for any $x \in X$ has
	\begin{equation*}
		P(x)=u(x,\cdots ,x).
	\end{equation*}
	then we call $P:X \to Y$ is a homogeneous polynomial of degree $k$.	
\end{defn}
Note that if $P$ is a homogeneous polynomial of degree $k$, then $P(\lambda X)=\lambda ^kP(x)$ for $x \in X $ and $\lambda  \in \mathbb{C}$.
\par
In this paper, the degree of a homogeneous polynomial will be indicated by a subscript. Namely, if $P_k$ is a homogeneous polynomial, then $P_k$ has the degree $k$. It's worth noting that if $P_k$ is a $k$-homogeneous polynomial from $X$ to $Y$, then exsits a unique symmetric $k-$linear mapping $u$ with
\begin{equation*}
	P(x)=u(x,\cdots ,x).
\end{equation*}
\par
Let $\mathcal{D} \in X$, $F$ is a homogeneous polynomial from $\mathcal{D}$ to $Y$. Then for any $z \in \mathcal{D} $, we denote $D^kF(z)$ is the $k$-Fr$\acute{e}$chet derivative of $F$ at $z$. If $\mathcal{D}$ contains the origin, then any holomorphic mapping $F:\mathcal{D} \to Y$ can be expanded into the series:
\begin{equation}
	F(z)=\sum_{n=0}^{\infty} \frac{1}{k!} D^kF(0)(z^k)
\end{equation}
Moreover, if $ \frac{1}{k!} D^kF(0)(z^k)$ is a$k$-homogeneous polynomial, then we use the notation $P_k(z)= \frac{1}{k!} D^kF(0)(z^k)$ alone this paper. Note that if $\mathcal{D}$ is a bounded balance domain in a complex Banach space X, and $F(\mathcal{D})$ is bounded, then for any $r \in \left ( 0,1 \right ) $ has(1.2) be converged in $r\mathcal{D}$.
\par
\begin{defn}
	Let $F$ is a holomorphic mapping from $\mathcal{D}$ to $Y$. For any $k \in \mathbb{N} $, if there has
	\begin{equation*}
		F(0)=0,DF(0)=0,\dots ,D^{k-1}F(0)=0 \quad but \quad D^kF(0)\ne0,
	\end{equation*}
	then we called $z=0$ is a zero of order $k$ of $F$.
\end{defn}
\begin{lem}\cite{ref25}
	Let $B_X$ and $B_Y$ are the unit balls of the Banach spaces $X$ and $Y$,respectively. We called the holomorphic mapping $v:B_X \to B_Y$ with $v(0)=0$ as Schwarz mapping.
	If $z=0$ is the zero of order k of $v$, then there has estimate holds
	\begin{equation}
		\left \|v(z) \right \|_Y \le \left \| z \right \| ^k_X,\quad z\in B_X.
	\end{equation}
\end{lem}
\begin{defn}
	Let $L\left ( X,\mathbb{C}  \right ) $ is a set of continuous linear operators from $X$ to $\mathbb{C}$. For any $x\in X\setminus \left \{ 0 \right \} $, let
	\begin{equation*}
		T(x)={l_x \in L\left ( X,\mathbb{C}  \right ):l_x(x)=\left \| x \right \|,\left \| l_x \right \|=1}.
	\end{equation*}
We can prove that this set is non-empty by using the Hahn-Banach theorem.
\end{defn}
The classical Schwarz lemma sates that for any holomorphis mappings $f$ for $\mathbb{D}$ to $\mathbb{D}$ with $f(0)=0$, then for all $z \in \mathbb{D}$ has $\left | f(z) \right | \le \left | z \right | $\cite{ref19}. Lindel$\ddot{o}$f removing the assumption that "the origin is a fixed point",then the classical Schwarz lemma (\cite{ref31}, Proposition 2.2.2) was improved. And based on Lindel$\ddot{o}$f's conclusion, Chen et al.\cite{ref11} obtained the Schwarz lemma between the unit spheres of complex Banach spaces as follows:
\begin{lem}\cite{ref11}
	Suppose that $B_X$ and $B_Y$ are the unit balls of complex Banach spaces $X$ and $Y$,respectively. Let $f:B_X \to B_Y$ be a holomorphic mapping. Then
	\begin{equation}
		\left \| f(z) \right \|_Y \le \frac{\left \| f(0) \right \|_Y+\left \| z \right \|_X }{1+\left \| f(0) \right \|_Y\left \| z \right\|_X }, for z \in B_X.
	\end{equation}
	This estimate is sharp with equality for each value of $\left \| f(0) \right \|_Y$ and for each $z \in B_X$.
\end{lem}
\par
\begin{lem}\cite{ref35}
Suppose that $f \in \mathcal{B}$ and $f(z) = \sum_{n=0}^{\infty} a_n z^n$. Then the following inequalities hold.
\begin{enumerate}
	\item[(a)] $|a_{2n+1}| \leq 1 - |a_0|^2 - \cdots - |a_n|^2,\ n = 0,1,\dots$,
	\item[(b)] $|a_{2n}| \leq 1 - |a_0|^2 - \cdots - |a_{n-1}|^2 - \frac{|a_n|^2}{1+|a_0|},\ n = 1,2,\dots$.
\end{enumerate}
Further, to have equality in (a) it is necessary that $f$ is a rational function of the form
	\begin{equation}
		f(z)=\frac{a_0+a_1z+\cdots +a_nz^n+\epsilon z^{2n}}{1+(\overline{a_n}z^n)+\cdots +\overline{a_0}z^{zn+1})\epsilon } ,\quad  \left| \epsilon \right|=1,
	\end{equation}
	and to have equality in (b) it is necessary that $f$ is a rational function of the form
	\begin{equation}
		f(z)=\frac{a_0+a_1z+\cdots +\frac{a_n}{1+\left | a_0 \right | } z^n+\epsilon z^{2n}}{1+(\frac{\overline{a_n}}{1+\left | a_0 \right | } z^n)+\cdots +\overline{a_0}z^{2n})\epsilon },\quad  \left| \epsilon \right|=1,
	\end{equation}
where the condition $a_0\overline{a_n}^2 \epsilon$ is non-positive real..
\end{lem}
By performing a simple calculation using Lemma 3, we can obtain
\begin{lem}\cite{ref35}
	Suppose that $f \in \mathcal{B}$ and $f(z) = \sum_{n=0}^{\infty} a_n z^n$. Then the inequality
	\begin{equation}
		M_f(r) \le \left | a_0 \right |+\frac{r}{1-r}(1- \left | a_0 \right |^2)-(\frac{1}{1+\left | a_0\right |}+\frac{r}{1-r})\left \| f_0 \right \|^2
	\end{equation}
	is valid for $r \in [0,1)$. Equality is attained for $f(z)=1,\ z\in\mathbb{D}$.
\end{lem}
\begin{rem}
	For (6) we consider that $a_0=0$ and obtain that
	\begin{equation*}
		\sum_{n=1}^{\infty}\left | a_n\right |r^n \le \frac{1}{1-r}(r-\|f_0\|^2).
	\end{equation*}
\end{rem}
\begin{lem}\cite{ref32}
	For any $f\in \mathcal{B} , r\in [0,1), t=\frac{N-1}{2}$, then there has the inequality
	\begin{align}
		S(r):=\sum_{n=N}^{\infty} \left | a_n \right | r^n
		&+\operatorname{sgn}(t)\sum_{n=1}^{t} \left | a_n \right |^2\frac{r^N}{1-r} +\left(\frac{1}{1+\left | a_0 \right | }+\frac{1}{1-r}\right)\sum_{n=t+1}^{\infty}\left |a_n  \right |^2r^{2n} \notag\\
		&\le (1-\left |a_0  \right |^2 )\frac{r^N}{1-r}
	\end{align}
holds.
\end{lem}
\begin{lem}\cite{ref10}
	Let $m\in \mathbb{N}$ and $p\in(0,2]$. For $a\in[0,1]$, consider
	\begin{equation*}
	a\mapsto D_{p,m}(a)=\left[\left(\frac{a+r^{m}}{1+ar^{m}}\right)^{p}-1\right]\varphi_{0}(r)+(1-a^{2})N(r),
\end{equation*}
	where $\varphi_{0}(r)$ and $N(r)$ are some nonnegative continuous functions defined on $[0,1)$. Also, suppose that
	\begin{equation*}
	\Psi_{p,m}(r)=p\left(\frac{1-r^{m}}{1+r^{m}}\right)\varphi_{0}(r)-2N(r)
\end{equation*}
	and $R:=R(m,p)$ is the minimal positive root in $(0,1)$ of the equation $\Psi_{p,m}(r)=0$. If $\Psi_{p,m}(r)\ge 0$ for $0\le r\le R$, then $D_{p,m}(a)\le 0$ for $0\le r\le R$.
\end{lem}
Let $\mathbb{B}_Y$ is the unit ball of a $\mathrm{JB}^*$-triple $Y$. Let $\operatorname{Aut}(\mathbb{B}_Y)$ denote the set of biholomorphic mappings of $\mathbb{B}_Y$ onto itself. Then $\operatorname{Aut}(\mathbb{B}_Y)$ is a transitive group.
\begin{lem}\cite{ref22}
	Let $F\colon \mathbb{U}\to \mathbb{B}_Y$ be a holomorphic mapping with
	\begin{equation*}
	F(\zeta) = a + \sum_{s=1}^{\infty} P_s(\zeta),\quad \zeta\in \mathbb{U},
    \end{equation*}
	where $P_s(\zeta) = \frac{1}{s!} D^s F(0)(\zeta^s)$. Assume that $a\in \mathbb{B}_Y$. Then, we have
	\begin{equation*}
	\frac{\|D\varphi_a(a)P_s(\zeta)\|}{\|D\varphi_a(a)\|} \le 1-\|a\|^2,\quad s\ge 1,\ \zeta\in \mathbb{U},
    \end{equation*}
	where $\varphi_a\in \operatorname{Aut}(\mathbb{B}_Y)$ such that $\varphi_a(a)=0$.
\end{lem}

\section{Improved Bohr-Rogosinski inequalities for holomorphic mappings with values in the closure of the unit polydisc of ${\mathbb{C}}^n $}\label{sec3}

In this section we have improved the inequalities established in Theorem E and Theorem F by using the $\left \| f_0 \right \| ^ 2$. First, we improve Theorem E and obtain the following theorem:
	\begin{thm}
	Let $B_X$ is a unit ball of the complex banach space $X$, $\mathbb{U}^n={z=(z_1,z_2,\dots,z_n)\in\mathbb{C}^n,\|z\|_\infty<1}$ is the unit polydics of $\mathbb{C}^n$, $a=(a_1,a_2,\dots,a_n)\in\mathbb{U}^n$, F is a holomorphic mapping from $B_X $ to $\overline {\mathbb{U}}^n$  with the form	
	\begin{align*}
		F(z)=a+\sum_{s=1}^{\infty} P_s(z),\quad z \in B_X,
	\end{align*}
	where $P_s(z)=\frac{1}{s!} D^sF(0)(z^s)$.
	Let $\left |a \right |=\| a \|_\infty$ holds for any $j=1,\dots ,n$, $\left \| F_0 \right \| ^2_\infty=\sum_{s=1}^{\infty}\left \| P_s(z) \right \|^2_\infty$. Then for any $\left \| z \right \|  =r \le R_1$, we have
	\begin{align}
		\left \| a \right \| _\infty +\sum_{n=1}^{\infty} \left \| P_s(z) \right \| _\infty+(\frac{1}{1+\left \| a \right \| _\infty} +\frac{r}{1-r} )\left \| F_0 \right \|^2_\infty \le 1.
	\end{align}
	The constant $R_1=\frac{2}{1+\left \| a \right \| _\infty}$ cannot be improved.
    \end{thm}
    \begin{proof}
	Let $z_0=\frac{z}{\left \| z \right \| } \in \partial \mathbb{U} ^n$ be fixed and $f_j(\xi )=F_j(\xi z_0),\xi \in \mathbb{U} $, where $F=(F_1,\dots ,F_n)$.\\
	Then $f_j(\xi )\in H(\mathbb{U},\mathbb{U} )$ and
	\begin{equation*}
		f_j(\xi )=a_j+\sum_{s=1}^{\infty} (P_s)_j(z_0)\xi ^s, \quad \xi \in \mathbb{U}.
	\end{equation*}
	Therefore, for each fixed j, we have
	\begin{equation*}
	(P_s)_j(z_0)\le 1-\left | a_j \right | ^2=1-\left \| a \right \| ^2_{\infty}.
	\end{equation*}
	This implies that
	\begin{equation}
		\left \| P_s(z_0) \right \| \le 1-\left \| a \right \| ^2_{\infty}.
	\end{equation}
	By the Lemma 4, we have
	\begin{align}
		\left | a_j \right | &+\sum_{s=1}^{\infty} \left | (P_s)_j(z_0)) \right | \left | \xi  \right | ^s \notag \\
		&\le \left | a_j \right | +(1-\left | a_j \right | ^2)\frac{\left | a_j \right | }{1-\left | a_j \right | } -(\frac{1}{1+\left | a_j \right | } +\frac{\left | \xi \right | }{1-\left | \xi \right |} )\left \| f_{j0} \right \| ^2 \notag \\
		&\le \left \| a \right \|_{\infty}  +(1-\left \| a \right \|_{\infty}  ^2)\frac{\left | \xi \right | }{1-\left | \xi \right | } -(\frac{1}{1+\left \| a \right \|_{\infty}  } +\frac{\left | \xi \right | }{1-\left | \xi \right |} )\sum_{s=1}^{\infty}\left \| P_s(z_0) \right \|_{\infty}^2\left | \xi \right |^{2s}
	\end{align}
	Let $x=\left \| a \right \|_{\infty}\in [0,1), r=\left \| z \right \|_{\infty}$, then $ r \in [0,1), z=rz_0$, by using estimates (9), (10), we have
	\begin{align*}
		\left \| a \right \|_{\infty}&+\sum_{s=1}^{\infty}\left \| P_s(z) \right \|_{\infty}+(\frac{1}{1+\left \| a \right \|_{\infty}} +\frac{r}{1-r}  )\left \| F_0 \right \|^2_{\infty} \\
		&\le \left \| a \right \|_{\infty}+(1-\left \| a \right \|_{\infty}^2)\frac{r}{1-r}\\
		&=1+\frac{(-x^2-x+2)r+(x-1)}{1-r}.
	\end{align*}
	Let $\Psi_x (r)=(-x^2-x+2)r+(x-1)$. The inequality holds $\le 1$ if and only if $\Psi_x (r) \le 0$. Because of
	\begin{equation*}
		\frac{\partial \Psi_x (r)}{\partial r} = -x^2-x+2 \ge 0, \quad x\in [0,1].
	\end{equation*}
	It means that $\Psi_x (r)$ is monotonically increasing on $r \in [0,1)$, so we have
	\begin{equation*}
		\Psi(r) \le \Psi(\frac{1}{2+\left \| a \right \|_{\infty} } ) =1.
	\end{equation*}
	Hence,we obtain $R_1=\frac{1}{2+\left \| a \right \|_{\infty} }$.
\par
	Next, we will show that the number $R_1=\frac{1}{2+\left \| a \right \|_{\infty} }$ is optimal.
	Let $z_0\in \partial B_X$ and $\mathbf{u}=(u_1,\dots,u_n)\in \partial \mathbb{U}^n$ with $|u_1|=\cdots=|u_n|=1$ be fixed. For $\lambda\in(0,1)$, let
	\begin{equation*}
	F(z)=f(l_{z_0}(z))\mathbf{u},\quad z\in B_X,
    \end{equation*}
	where
	\begin{equation*}
	f(\zeta)=\frac{\lambda-\zeta}{1-\lambda\zeta},\quad \zeta\in \mathbb{U}
    \end{equation*}
	and $l_{z_0}\in T(z_0)$. Then we have
	\begin{align*}
		\left \| a \right \| _\infty &+\sum_{n=1}^{\infty} \left \| P_s(z) \right \| _\infty+(\frac{1}{1+\left \| a \right \| _\infty} +\frac{r}{1-r} )\left \| F \right \|^2_\infty\\
		&=\lambda+(1-\lambda^2)\frac{r}{1-\lambda}+(\frac{1}{1+\lambda} +\frac{r}{1-r} )(1-\lambda^2)^2\frac{r^2}{1-\lambda^2r^2} \\
		&=1+\frac{(1-\lambda)[(2+\lambda)r-1]}{1-r}
	\end{align*}
	Obviously, the inequality will $>$ 1 when $r > \frac{1}{2+\lambda} $.Thus, $R_1$ is optimal. This completes the proof.
\end{proof}

We now already know that through replacing $\mathit{f}(z)$ with $\mathit{f}(0)$, the Bohr inequality can be optimized. We replaced $\left \| a \right \| _\infty $ with $\left \| F(v(z)) \right \| _\infty^p$ for the (8). And make the same method improvement to the Bohr-Rogosinski type inequality, we can obtain the following theorem:
\begin{thm}
	Let $B_X$ is a unit ball of the complex banach space $X$, $\mathbb{U}^n={z=(z_1,z_2,\dots,z_n)\in\mathbb{C}^n,\|z\|_\infty<1}$ is the unit polydics of $\mathbb{C}^n$, $a=(a_1,a_2,\dots,a_n)\in\mathbb{U}^n$, F is a holomorphic mapping from $B_X $ to $\overline {\mathbb{U}}^n$  with the form	
	\begin{align*}
		F(z)=a+\sum_{s=1}^{\infty} P_s(z),\quad z \in B_X,
	\end{align*}
	where $P_s(z)=\frac{1}{s!} D^sF(0)(z^s)$.
	Let $\left |a_j \right |=\| a \|_\infty$ holds for any $j=1,\dots ,n$, $\left \| F_0 \right \| ^2_\infty=\sum_{s=1}^{\infty}\left \| P_s(z) \right \|^2_\infty$. Then for any $m \in \mathbb{N} , p \in (0,2]$ and $\left \| z \right \| _\infty=r \le R_2$, we have
	\begin{align}
		\left \| F(v(z)) \right \| _\infty^p +\sum_{n=1}^{\infty} \left \| P_s(z) \right \| _\infty+(\frac{1}{1+\left \| a \right \| _\infty} +\frac{r}{1-r} )\left \| F_0 \right \|^2_\infty \le 1.
	\end{align}
	Where $v(z): B_X \to B_X$ is a Schwarz mapping having $z=0$ as a zero of order m, and constant $R_2$ is the unique root in $(0,1)$ of the equation
	\begin{equation}
	p(\frac{1-r^m}{1+r^m} )-2\frac{r}{1-r}=0
	\end{equation}
	The constant $R_2$ cannot be improved.
\end{thm}
\begin{proof}
	Let $z_0=\frac{z}{\left \| z \right \| } \in \partial \mathbb{U} ^n$ be fixed and $f_j(\xi )=F_j(\xi z_0),\xi \in \mathbb{U} $, where $F=(F_1,\dots ,F_n)$.\\
	Then $f_j(\xi )\in H(\mathbb{U},\mathbb{U} )$ and
	\begin{equation*}
		f_j(\xi )=a_j+\sum_{s=1}^{\infty} (P_s)_j(z_0)\xi ^s, \quad \xi \in \mathbb{U}.
	\end{equation*}
	Therefore, for each fixed j, we have
	\begin{equation*}
		(P_s)_j(z_0)\le 1-\left | a_j \right | ^2=1-\left \| a \right \| ^2_{\infty}.
	\end{equation*}
	This implies that
	\begin{equation}
		\left \| P_s(z_0) \right \| \le 1-\left \| a \right \| ^2_{\infty}.
	\end{equation}
		Let $x=\left \| a \right \|_{\infty}\in [0,1), r=\left \| z \right \|_{\infty}$, then $ r \in [0,1), z=rz_0$, by using Lemma 2,5 and estimates (13), we have
	\begin{align*}
		\left \| F(v(z)) \right \|_{\infty}^p &+\sum_{s=1}^{\infty}\left \| P_s(rz_0) \right \|_{\infty}+(\frac{1}{1+\left \| a \right \|_{\infty}} +\frac{r}{1-r}  )\left \| F_0 \right \|^2_{\infty} \\
		&\le \left \| F(v(z)) \right \|_{\infty}^p+(1-\left \| a \right \|_{\infty}^2)\frac{r}{1-r}\\
		&\le (\frac{\left \| F(0) \right \|_\infty+\left \| v(rz_0) \right \|_\infty }{1+\left \| F(0) \right \|_\infty\left \| v(rz_0) \right \|_\infty})^p+(1-x^2)\frac{r}{1-r}\\
		&\le (\frac{x+\left \| rz_0 \right \|_\infty^m }{1+x\left \| rz_0 \right \|_\infty^m})^p+(1-x^2)\frac{r}{1-r}\\
		&=1+[(\frac{x+r^m }{1+xr^m})^p+(1-x^2)\frac{r}{1-r}]\\
		&=1+\Phi _{m,\mathbb{N}}^p(r),
	\end{align*}
	which is less than or equal 1 provided $\Phi _{m,\mathbb{N}}^p(r)$, where
	\begin{equation}
	\Phi _{m,\mathbb{N}}^p(r)=(\frac{x+r^m }{1+xr^m})^p+(1-x^2)\frac{r}{1-r}-1.
	\end{equation}
	Hence, by Lemma 6, let $x=a, \varphi _0(r)\equiv 1, N(r)=\frac{r}{1-r} $, we can obtain $\Psi _{m,\mathbb{N}}^p \le 0$ for $r \le R_2$, where $R_2$ is the unique root in (0,1) of the equation (12).
	\par
	Next, we will show that the number $R_2$ is optimal.
	Let $z_0\in \partial B_X$ and $\mathbf{u}=(u_1,\dots,u_n)\in \partial \mathbb{U}^n$ with $|u_1|=\cdots=|u_n|=1$ be fixed. For $\lambda\in(0,1)$, let
	\begin{equation*}
		F(z)=f(l_{z_0}(z))\mathbf{u},\quad z\in B_X,
	\end{equation*}
	where
	\begin{equation*}
		f(\zeta)=\frac{\lambda-\zeta}{1-\lambda\zeta},\quad \zeta\in \mathbb{U}
	\end{equation*}
	and $l_{z_0}\in T(z_0)$.
	Let $v(z)=-l_{z_0}(z)^{m-1}z$, then we have
	\begin{align*}
		\left \| F(v(z)) \right \|_{\infty}^p &+\sum_{s=1}^{\infty}\left \| P_s(rz_0) \right \|_{\infty}+(\frac{1}{1+\left \| a \right \|_{\infty}} +\frac{r}{1-r}  )\left \| F_0 \right \|^2_{\infty} \\
		&=(\frac{\lambda+r^m }{1+\lambda r^m})^p+(1-\lambda^2)\frac{r}{1-\lambda r}+(\frac{1}{1+\lambda }\\ &+\frac{r}{1-r})(1-\lambda^2)^2\frac{r^2}{1-\lambda^2r^2}\\
		&=1+[(\frac{\lambda+r^m }{1+\lambda r^m})^p+(1-\lambda^2)\frac{r}{1-r}]\\
		&=1+(1-\lambda)\Psi  _{m,\mathbb{N}}^p(\lambda,r)
	\end{align*}
	where
	\begin{equation*}
	\Psi  _{m,\mathbb{N}}^p(\lambda,r)=\frac{1}{1-\lambda}[(\frac{\lambda+r^m }{1+\lambda r^m})^p]-1+(1+\lambda)\frac{r}{1-r}
	\end{equation*}
	Let $r \in (R_2, 1)$ be arbitrarily fixed. Since $R_2$ satisfies the equation (12), it follows that
	\begin{equation*}
	-p\left(\frac{1-r^m}{1+r^m}\right) + \frac{2r^N}{1-r} > 0
	\end{equation*}
	for $r \in (R_2, 1)$. Therefore
	\begin{equation*}
	\lim_{\lambda\to 1^-} \Psi_{m,N}^p(\lambda,r)
	= -p\left(\frac{1-r^m}{1+r^m}\right) + \frac{2r^N}{1-r}
	> 0,
    \end{equation*}
	and hence there exists $\lambda \in (0,1)$ such that $\Psi_{m,N}^p(\lambda,r) > 0$. This implies that
	\begin{equation*}
	\left \| F(v(z)) \right \|_{\infty}^p+\sum_{s=1}^{\infty}\left \| P_s(rz_0) \right \|_{\infty}+(\frac{1}{1+\left \| a \right \|_{\infty}} +\frac{r}{1-r}  )\left \| F_0 \right \|^2_{\infty} > 1.
    \end{equation*}
	Thus, $R_2$ is optimal. This completes the proof.
\end{proof}

\section{Improved Bohr-Rogosinski inequalities for holomorphic mappings with values in the closure of the unit ball of a $JB*-$triple}\label{sec4}
In this section, we will generalize the Bohr type inequality (8) and Bohr-Rogosinski type radius (11) to the case of holomorphic mappings on the unit ball of a complex Banach space with values in $\mathbb{B}_Y$.

\begin{thm}
	Let $B_X$ is a unit ball of the complex banach space $X$, $\mathbb{B}_Y$ is the unit ball of a $\mathrm{JB}^*$-triple $Y$. Let F is a holomorphic mapping from $B_X $ to $\mathbb{B}_Y$ with the form	
	\begin{align*}
		F(z)=a+\sum_{s=1}^{\infty} P_s(z),\quad z \in B_X,
	\end{align*}
	where $P_s(z)=\frac{1}{s!} D^sF(0)(z^s)$.
	Assume that $a \in \mathbb{B}_Y$. Then for $r \le R_1$, we have
	\begin{align}
		\left \| a \right \| _\infty +\sum_{n=1}^{\infty} \frac{\|D\varphi_a(a)P_s(\zeta)\|}{\|D\varphi_a(a)\|}+(\frac{1}{1+\left \| a \right \| _\infty} +\frac{r}{1-r} )\left \| F_0 \right \|^2_\infty \le 1.
	\end{align}
	The constant $R_1=\frac{2}{1+\left \| a \right \| _\infty}$ cannot be improved.
\end{thm}
\begin{proof}
	Let $z_0=\frac{z}{\left \| z \right \| } \in \partial \mathbb{U} ^n$ be fixed and $f_j(\xi )=F_j(\xi z_0), \xi \in \mathbb{U} $, where $F=(F_1,\dots ,F_n)$.\\
	Then $f_j(\xi )\in H(\mathbb{U},\overline{\mathbb{B}_Y} )$ and
	\begin{equation*}
		f_j(\xi )=a_j+\sum_{s=1}^{\infty} (P_s)_j(z_0)\xi ^s, \quad \xi \in \mathbb{U}.
	\end{equation*}
	Let $a=f(0)\in \mathbb{B}_Y$, by Lemma 7,we have
	\begin{equation}
	\frac{\|D\varphi_a(a)P_s(\zeta)\|}{\|D\varphi_a(a)\|}\le 1-\|a\|_\infty ^2, \quad s\ge 1.
	\end{equation}
	Let $x=\left \| a \right \|_{\infty}\in [0,1), r=\left \| z \right \|_{\infty}$, then $ r \in [0,1), z=rz_0$, by using Lemma 1 and estimates (16), we have
	\begin{align*}
		\left \| a \right \|_{\infty}&+\sum_{s=1}^{\infty}	\frac{\|D\varphi_a(a)P_s(\zeta)\|}{\|D\varphi_a(a)\|}+(\frac{1}{1+\left \| a \right \|_{\infty}} +\frac{r}{1-r}  )\left \| F_0 \right \|^2_{\infty} \\
		&\le \left \| a \right \|_{\infty}+(1-\left \| a \right \|_{\infty}^2)\frac{r}{1-r}\\
		&=1+\frac{(-x^2-x+2)r+(x-1)}{1-r}.
	\end{align*}
	Let $\Psi_x (r)=(-x^2-x+2)r+(x-1)$. The inequality holds $\le 1$ if and only if $\Psi_x (r) \le 0$. Because of
	\begin{equation*}\label{eq-4.1}
		\frac{\partial \Psi_x (r)}{\partial r} = -x^2-x+2 \ge 0, \quad x\in [0,1].
	\end{equation*}
	It means that $\Psi_x (r)$ is monotonically increasing on $r \in [0,1)$, so we have
	\begin{equation*}
		\Psi(r) \le \Psi(\frac{1}{2+\left \| a \right \|_{\infty} } ) =1.
	\end{equation*}
	Hence,we obtain $R_1=\frac{1}{2+\left \| a \right \|_{\infty} }$.
	\par
	Next, we will show that the number $R_1=\frac{1}{2+\left \| a \right \|_{\infty} }$ is optimal.
	Let $z_0\in \partial B_X$ and $\mathbf{u}=(u_1,\dots,u_n)\in \partial \mathbb{B}_Y$ be fixed. For $\lambda\in(0,1)$, let
	\begin{equation*}
		F(z)=f(l_{z_0}(z))\mathbf{u},\quad z\in B_X,
	\end{equation*}
	where
	\begin{equation*}
		f(\zeta)=\frac{\lambda-\zeta}{1-\lambda\zeta},\quad \zeta\in \mathbb{U}
	\end{equation*}
	and $l_{z_0}\in T(z_0)$. Then we have
	\begin{align*}
		\left \| a \right \| _\infty &+\sum_{n=1}^{\infty} \frac{\|D\varphi_a(a)P_s(\zeta)\|}{\|D\varphi_a(a)\|}+(\frac{1}{1+\left \| a \right \| _\infty} +\frac{r}{1-r} )\left \| F \right \|^2_\infty\\
		&=\lambda+(1-\lambda^2)\frac{r}{1-\lambda}+(\frac{1}{1+\lambda} +\frac{r}{1-r} )(1-\lambda^2)^2\frac{r^2}{1-\lambda^2r^2} \\
		&=1+\frac{(1-\lambda)[(2+\lambda)r-1]}{1-r}
	\end{align*}
	Obviously, the inequality will over 1 when $r > \frac{1}{2+\lambda} $.Thus, $R_1$ is optimal. This completes the proof.
\end{proof}

\begin{thm}
	Let $B_X$ is a unit ball of the complex banach space $X$, $\mathbb{B}_Y$ is the unit ball of a $\mathrm{JB}^*$-triple $Y$. Let F is a holomorphic mapping from $B_X $ to $\mathbb{B}_Y$ with the form	
	\begin{align*}
		F(z)=a+\sum_{s=1}^{\infty} P_s(z),\quad z \in B_X,
	\end{align*}
	where $P_s(z)=\frac{1}{s!} D^sF(0)(z^s)$.
    Assume that $a \in \mathbb{B}_Y$.Then for any $m \in \mathbb{N} , p \in (0,2]$ and $\left \| z \right \| _\infty=r \le R_2$, we have
	\begin{align}
		\left \| F(v(z)) \right \| _\infty^p ++\sum_{s=1}^{\infty}	\frac{\|D\varphi_a(a)P_s(\zeta)\|}{\|D\varphi_a(a)\|}+(\frac{1}{1+\left \| a \right \|_{\infty}} +\frac{r}{1-r}  )\left \| F_0 \right \|^2_{\infty} \le 1.
	\end{align}
	Where $v(z): B_X \to B_X$ is a Schwarz mapping having $z=0$ as a zero of order $m$, $\varphi_a \in \operatorname{Aut}(\mathbb{B}_Y)$ such that $\varphi_a(a)=0$ and $R_2$ is the unique root in $(0,1)$ of the equation (12). The number $R_2$ cannot be improved.
    \end{thm}
    \begin{proof}
	Let $z_0=\frac{z}{\left \| z \right \| } \in \partial \mathbb{U} ^n$ be fixed and $f_j(\xi )=F_j(\xi z_0), \xi \in \mathbb{U} $, where $F=(F_1,\dots ,F_n)$.\\
    Then $f_j(\xi )\in H(\mathbb{U},\overline{\mathbb{B}_Y} )$ and
    \begin{equation*}
	f_j(\xi )=a_j+\sum_{s=1}^{\infty} (P_s)_j(z_0)\xi ^s, \quad \xi \in \mathbb{U}.
    \end{equation*}
    Let $a=f(0)\in \mathbb{B}_Y$,$x=\left \| a \right \|_{\infty}\in [0,1), r=\left \| z \right \|_{\infty}$, then $ r \in [0,1), z=rz_0$,  by using Lemma 2,5 and estimates (13), (16), we have
	\begin{align*}
		\left \| F(v(z)) \right \|_{\infty}^p &++\sum_{s=1}^{\infty}	\frac{\|D\varphi_a(a)P_s(\zeta)\|}{\|D\varphi_a(a)\|}+(\frac{1}{1+\left \| a \right \|_{\infty}} +\frac{r}{1-r}  )\left \| F_0 \right \|^2_{\infty}  \\
		&\le \left \| F(v(z)) \right \|_{\infty}^p+(1-\left \| a \right \|_{\infty}^2)\frac{r}{1-r}\\
		&\le (\frac{\left \| F(0) \right \|_\infty+\left \| v(rz_0) \right \|_\infty }{1+\left \| F(0) \right \|_\infty\left \| v(rz_0) \right \|_\infty})^p+(1-x^2)\frac{r}{1-r}\\
		&\le (\frac{x+\left \| rz_0 \right \|_\infty^m }{1+x\left \| rz_0 \right \|_\infty^m})^p+(1-x^2)\frac{r}{1-r}\\
		&=(\frac{x+r^m }{1+xr^m})^p+(1-x^2)\frac{r}{1-r}
	\end{align*}
    As in the proof of Theorem 2, we obtain
	\begin{equation*}
		(\frac{x+r^m }{1+xr^m})^p+(1-x^2)\frac{r}{1-r}\le 1.
	\end{equation*}
    for $r \le R_2$, where $R_2$ is the unique root in (0,1) of the equation (12).
	\par
	Next, we will show that the number $R_2$ is optimal.
	Let $z_0\in \partial B_X$ and $\mathbf{u}=(u_1,\dots,u_n)\in \partial \mathbb{B}_Y$ be fixed. For $\lambda\in(0,1)$, let
	\begin{equation*}
		F(z)=f(l_{z_0}(z))\mathbf{u},\quad z\in B_X,
	\end{equation*}
	where
	\begin{equation*}
		f(\zeta)=\frac{\lambda-\zeta}{1-\lambda\zeta},\quad \zeta\in \mathbb{U}
	\end{equation*}
	and $l_{z_0}\in T(z_0)$.
	Let $v(z)=-l_{z_0}(z)^{m-1}z$, then we have
	\begin{align*}
		\left \| F(v(z)) \right \|_{\infty}^p &+\sum_{s=1}^{\infty}\left \| P_s(rz_0) \right \|_{\infty}+(\frac{1}{1+\left \| a \right \|_{\infty}} +\frac{r}{1-r}  )\left \| F_0 \right \|^2_{\infty} \\
		&=(\frac{\lambda+r^m }{1+\lambda r^m})^p+(1-\lambda^2)\frac{r}{1-\lambda r}+(\frac{1}{1+\lambda }\\ &+\frac{r}{1-r})(1-\lambda^2)^2\frac{r^2}{1-\lambda^2r^2}\\
		&=1+[(\frac{\lambda+r^m }{1+\lambda r^m})^p+(1-\lambda^2)\frac{r}{1-r}]\\
	\end{align*}
	As in the proof of Theorem 2, we obtain for $r \in (R_2, 1)$, there exists $\lambda \in (0,1)$ such that $\Psi_{m,N}^p(\lambda,r) > 0$. This implies that
	\begin{equation*}
	\left \| F(v(z)) \right \|_{\infty}^p+\sum_{s=1}^{\infty}\left \| P_s(rz_0) \right \|_{\infty}+(\frac{1}{1+\left \| a \right \|_{\infty}} +\frac{r}{1-r}  )\left \| F_0 \right \|^2_{\infty} > 1.
	\end{equation*}
	Thus, $R_2$ is optimal. This completes the proof.
\end{proof}
% ------------------------------------------------------------------------


\subsection*{Data Availability Statement}Not applicable. The manuscript has no associated data.

\subsection*{Competing Interests}The authors have no relevant financial or non-financial interests to disclose.

\subsection*{Author Contributions}All authors contributed to the study conception and design.
The first draft of the manuscript was written by Liangpeng Xiong and all authors commented on previous versions of the manuscript. All authors read and approved the final manuscript.

\begin{thebibliography}{40}
	
	\bibitem{ref1}Abu-Muhanna Y, Ali R M, Ng Z C, Hasni S F M.: Bohr radius for subordinating families of analytic functions and bounded harmonic mappings. Journal of Mathematical Analysis and Applications, \textbf{420}, 124--136 (2014)
	
	\bibitem{ref2}Aizenberg L.: Generalization of results about the Bohr radius for power series. Studia Mathematica, \textbf{180}, 161--168 (2007)
	
	\bibitem{ref3}Aizenberg L.: Multidimensional analogues of Bohr's theorem on power series. Proceedings of the American Mathematical Society, \textbf{128}(4), 1147--1155 (2000)
	
	\bibitem{ref4}Aizenberg L, Aytuna A, Djakov P.: Generalization of a theorem of Bohr for bases in spaces of holomorphic functions of several complex variables. Journal of Mathematical Analysis and Applications, \textbf{258}(1), 429--447 (2001)
	
	\bibitem{ref5}Alkhaleefah S A, Kayumov I R, Ponnusamy S.: On the Bohr inequality with a fixed zero coefficient. Proceedings of the American Mathematical Society, \textbf{147}, 5263--5274 (2019)
	
	\bibitem{ref6}Bayart F, Pellegrino D, Seoane-Sepúlveda J B.: The Bohr radius of the n-dimensional polydisk is equivalent to $(\log n)/n$. Advances in Mathematics, \textbf{264}, 726--746 (2014)
	
	\bibitem{ref7}Bénéteau C, Dahlner A, Khavinson D.: Remarks on the Bohr phenomenon. Computational Methods and Function Theory, \textbf{4}(1), 1--19 (2004)
	
	\bibitem{ref8}Boas H P, Khavinson D.: Bohr's power series theorem in several variables. Proceedings of the American Mathematical Society, \textbf{125}(10), 2975--2979 (1997)
	
	\bibitem{ref9}Bohr H.: A theorem concerning power series. Proceedings of the London Mathematical Society, s2-\textbf{13}, 1--5 (1914)
	
	\bibitem{ref10}Chen K, Liu M S, Ponnusamy S.: Bohr-type inequalities for unimodular bounded analytic functions. Results in Mathematics, \textbf{78}(5), Paper No.183, 16 pp (2023)
	
	\bibitem{ref11}Chen S L, Hamada H, Ponnusamy S, Vijayakumar R.: Schwarz type lemmas and their applications in Banach spaces. Journal d'Analyse Mathématique, \textbf{152}, 181--216 (2024)
	
	\bibitem{ref12}Chu C H.: Jordan structures in geometry and analysis. Cambridge University Press, Cambridge (2012). (Cambridge Tracts in Mathematics 190)
	
	\bibitem{ref13}Chu C H.: Bounded symmetric domains in Banach spaces. World Scientific Publishing Co. Pte. Ltd., Hackensack, NJ (2021)
	
	\bibitem{ref14}Defant A, Frerick L.: A logarithmic lower bound for multi-dimensional Bohr radii. Israel Journal of Mathematics, \textbf{152}, 17--28 (2006)
	
	\bibitem{ref15}Defant A, Frerick L.: The Bohr radius of the unit ball of $\ell_p^n$. Journal für die reine und angewandte Mathematik, \textbf{660}, 131--147 (2011)
	
	\bibitem{ref16}Defant A, Frerick L, Ortega-Cerdà J, Ounaïes M, Seip K.: The Bohnenblust–Hille inequality for homogeneous polynomials is hypercontractive. Annals of Mathematics, \textbf{174}(1), 485--497 (2011)
	
	\bibitem{ref17}Dineen S, Timoney R M.: Absolute bases, tensor products and a theorem of Bohr. Studia Mathematica, \textbf{94}(3), 227--234 (1989)
	
	\bibitem{ref18}Djakov P B, Ramanujan M S.: A remark on Bohr’s theorem and its generalizations. Journal of Analysis, \textbf{8}, 65--77 (2000)
	
	\bibitem{ref19}Graham I, Kohr G.: Geometric Function Theory in One and Higher Dimensions. Marcel Dekker Inc., New York (2003)
	
	\bibitem{ref20}Hamada H.: Bohr phenomenon for analytic functions subordinate to starlike or convex functions. Journal of Mathematical Analysis and Applications, \textbf{499}, 125019 (2021)
	
	\bibitem{ref21}Hamada H, Honda T.: Some generalizations of Bohr's theorem. Mathematical Methods in the Applied Sciences, \textbf{35}(17), 2031--2035 (2012)
	
	\bibitem{ref22}Hamada H, Honda T.: Bohr phenomena for holomorphic mappings with values in several complex variables. Results in Mathematics, \textbf{79}, Article 239 (2024)
	
	\bibitem{ref23}Hamada H, Honda T.: Bohr radius for pluriharmonic mappings in separable complex Hilbert spaces. Bulletin of the Malaysian Mathematical Sciences Society, \textbf{47}, Article 47 (2024)
	
	\bibitem{ref24}Hamada H, Honda T, Kohr G.: Bohr's theorem for holomorphic mappings with values in homogeneous balls. Israel Journal of Mathematics, \textbf{173}, 177--187 (2009)
	
	\bibitem{ref25}Hamada H, Honda T, Kohr M.: Bohr–Rogosinski radius for holomorphic mappings with values in higher dimensional complex Banach spaces. Analysis and Mathematical Physics, \textbf{15}, Article 64 (2025)
	
	\bibitem{ref26}Hamada H, Honda T, Mizota Y.: Bohr phenomenon on the unit ball of a complex Banach space. Mathematical Inequalities \& Applications, \textbf{23}(4), 1325--1341 (2020)
	
	\bibitem{ref27}Huang Y, Liu M S, Ponnusamy S.: Refined Bohr-type inequalities with area measure for bounded analytic functions. Analysis and Mathematical Physics, \textbf{10}, 50 (2020)
	
	\bibitem{ref28}Kayumov I R, Khammatova D M, Ponnusamy S.: Bohr-Rogosinski phenomenon for analytic functions and Cesàro operators. Journal of Mathematical Analysis and Applications, \textbf{496}(2), 124824, 17 pp (2021)
	
	\bibitem{ref29}Kayumov I R, Ponnusamy S.: On a powered Bohr inequality. Annales Academiæ Scientiarum Fennicæ Mathematica, \textbf{44}, 301--310 (2019)
	
	\bibitem{ref30}Kaup W.: A Riemann mapping theorem for bounded symmetric domains in complex Banach spaces. Mathematische Zeitschrift, \textbf{183}, 503--529 (1983)
	
	\bibitem{ref31}Krantz S G.: Geometric function theory: Explorations in complex analysis. Birkhäuser, Boston (2006)
	
	\bibitem{ref32}Liu G, Liu Z, Ponnusamy S.: Refined Bohr inequality for bounded analytic functions. Bulletin des Sciences Mathématiques, \textbf{173}, 103054 (2021)
	
	\bibitem{ref33}Liu M S, Ponnusamy S.: Multidimensional analogues of refined Bohr’s inequality. Proceedings of the American Mathematical Society, \textbf{149}(5), 2133--2146 (2021)
	
	\bibitem{ref34}Paulsen V I, Popescu G, Singh D.: On Bohr's inequality. Proceedings of the London Mathematical Society, \textbf{85}, 493--512 (2002)
	
	\bibitem{ref35}Ponnusamy S, Vijayakumar R, Wirths K-J.: New inequalities for the coefficients of unimodular bounded functions. Results in Mathematics, \textbf{75}, 107 (2020)
	
	\bibitem{ref36}Ponnusamy S, Vijayakumar R, Wirths K J.: Improved Bohr’s phenomenon in quasi-subordination classes. Journal of Mathematical Analysis and Applications, \textbf{506}(1), Paper No.125645, 10 pp (2022)
	
	\bibitem{ref37}Rogosinski W.: Über Bildschranken bei Potenzreihen und ihren Abschnitten. Mathematische Zeitschrift, \textbf{17}, 260--276 (1923)
	
\end{thebibliography}


% ------------------------------------------------------------------------
\end{document}
% ------------------------------------------------------------------------
